\documentclass[10pt,a4paper]{article}

\usepackage[T1]{fontenc}
\usepackage[utf8]{inputenc}
\usepackage{lmodern}
\usepackage{amsmath,amssymb,amsthm}
\usepackage{booktabs,array}
\usepackage{geometry}
\usepackage{graphicx}
\usepackage{xcolor}
\usepackage{enumitem}
\usepackage{microtype}
\usepackage{hyperref}
\usepackage{fancyhdr}

\geometry{margin=1in}
\hypersetup{
  colorlinks=true,
  linkcolor=blue!55!black,
  citecolor=blue!55!black,
  urlcolor=blue!55!black,
  pdftitle={A Universal Zero-Witness Forgery against Miraidon-S},
  pdfauthor={ePrint Signature Audit}
}
\setlist{nosep,leftmargin=*}
\setlength{\parindent}{1em}
\setlength{\parskip}{0.25em}
\pagestyle{fancy}
\fancyhf{}
\fancyhead[L]{Zero-witness forgery against Miraidon-S}
\fancyhead[R]{7 October 2026}
\fancyfoot[C]{\thepage}
\renewcommand{\headrulewidth}{0.3pt}

\newtheorem{theorem}{Theorem}
\newtheorem{observation}{Observation}
\newcommand{\F}{\mathbb{F}}
\newcommand{\Com}{\mathsf{Com}}
\newcommand{\FSone}{\mathsf{FS}_1}
\newcommand{\FStwo}{\mathsf{FS}_2}
\newcommand{\rank}{\operatorname{rank}}
\newcommand{\GL}{\operatorname{GL}}
\newcommand{\matzero}{\mathbf{0}}
\newcommand{\vectzero}{\boldsymbol{0}}
\newcommand{\accept}{\mathsf{accept}}

\title{\bfseries A Universal Zero-Witness Forgery against Miraidon-S\\
\large Cryptanalysis of \emph{Miraidon: MinRank Identification}}
\author{ePrint Signature Audit\\
\small AI-assisted cryptanalysis report}
\date{7 October 2026}

\begin{document}
\maketitle

\begin{abstract}
We give a public-key-only universal fresh-message forgery against Miraidon-S
as specified in \emph{Miraidon: MinRank Identification}, IACR ePrint
2026/997, version~4.  The protocol is intended to prove knowledge of a
coefficient vector whose public matrix combination has rank exactly $r$.
The verifier, however, accepts a factorization $Q=LR$ without checking that
$Q$ has exact rank $r$, while the paper's printed homogeneous MinRank problem
also omits a nonzero condition on the coefficient vector.  A forger sets the
purported coefficient vector and all matrix responses to zero.  Both verifier
branches then hash the same zero matrix for every first challenge, eliminating
the fixed-weight Fiat--Shamir challenge-guessing cost.  The attack uses no
secret information, signing query, MinRank computation, or random-oracle
grinding, and succeeds for every public key and message.

We reproduce the complete Fiat--Shamir forgery at the Level-I, Level-III, and
Level-V parameter rows.  A full Level-I run using an exact Algorithm-5 key
took 0.49 seconds for the forgery and 4.94 seconds including key generation,
honest signing, and controls.  A separate implementation reproduced two
fresh forgeries at every parameter row.  An exact-rank check rejects every
forgery and preserves honest signatures.  The result was AI-discovered and
independently AI-reproduced; independent human verification is pending.
No originality claim is made for the general technique of detecting a
missing nondegeneracy or relation-membership check.
\end{abstract}

\noindent\fcolorbox{black!45}{black!4}{\parbox{0.94\textwidth}{
\textbf{Result status.} Practical full-parameter design attack: universal
fresh-message forgery with zero signing queries.  Discovery, cryptanalysis,
implementation, adversarial review, and independent reproduction were
AI-assisted.  Independent human verification is pending.}}

\section{Introduction}

Miraidon-S is a Fiat--Shamir signature derived from a five-pass MinRank
identification protocol.  Its target paper, by Cartor and Slaughter, is titled
\emph{Miraidon: MinRank Identification}~\cite{miraidon}.  The paper proposes
parameter sets for NIST security Levels I, III, and V and analyzes both
generic MinRank attacks and fixed-weight attacks on repeated five-pass
protocols.  The attack in this paper bypasses both analyses: a malformed
rank-zero response is accepted by both protocol branches, so the forger need
not predict either challenge.

Our contributions are as follows.
\begin{itemize}
  \item We identify the mismatch between the exact-rank private relation and
  the verifier's at-most-$r$ factorization check.  The printed homogeneous
  MinRank problem compounds the issue by allowing the zero coefficient
  vector.
  \item We derive an all-zero transcript that answers both second challenges
  for every nonzero first challenge.  It yields a universal public-key-only
  forgery after the same seed-tree, Merkle-tree, and Fiat--Shamir processing
  used by an honest signature.
  \item We validate the attack on all three published parameter rows using
  independent implementations, an exact Algorithm-5 Level-I key, honest and
  changed-message controls, an exact-rank-hardened verifier, and the authors'
  pinned interactive proof-of-concept.
  \item We give a concrete repair: check $\rank(Q)=r$ in every opened
  branch-zero response (or equivalently require both factors to have rank
  $r$), state a nondegenerate witness relation, correct the aggregate checks
  in Algorithm~6, and revise the soundness proof.
\end{itemize}

\paragraph{Attack-history statement.}
We located no previous public attack against exact Miraidon-S in searches as
of 7 October 2026.  The search covered the exact title, scheme name, ePrint
identifier, authors, repository records, and terms concerning zero witnesses,
rank zero, missing rank checks, soundness, and forgery.  This negative result
has moderate confidence: the target is recent, indexes may be incomplete, and
unpublished or unindexed observations cannot be excluded.  We therefore make
no absolute priority claim.

\paragraph{Technique attribution.}
Generic algebraic and combinatorial MinRank algorithms are well established;
examples include support-minors and related algebraic methods
\cite{bardet2020,bardet2023}.  Challenge-guessing attacks on Fiat--Shamir
signatures from five-pass identification schemes were developed by Kales and
Zaverucha~\cite{kales2020} and refined for fixed-weight repetition by
Battagliola et al.~\cite{battagliola2025revision,battagliola2025fixed}.
Miraidon's claimed parameters account for those costs.  Our attack uses
neither class of method.  Its essential observation is a missing
nondegeneracy check, a standard protocol-review issue for which we claim no
technique novelty.

\section{Target, notation, and claimed security}

\subsection{Exact target}

The target is Miraidon-S in IACR ePrint 2026/997, version~4, dated 2 June
2026~\cite{miraidon}.  The audited PDF has SHA-256 digest
\begin{center}
\small\ttfamily
88f0bcfb4f4a4e7079236d48369d5d692d0c39a89447b2bbe3d7bcd4dcfe21f4.
\end{center}
The accompanying proof-of-concept at
\url{https://github.com/FreemanSlaughter/Miraidon} is pinned to Git commit
\begin{center}
\small\ttfamily 289beaa76a75e8500faa6fa75e9c76796fe9f633.
\end{center}
The pinned \path{Miraidon_ZKP.py} has SHA-256 digest
\begin{center}
\small\ttfamily 0523b0b89d3be7cca86b23b933e18a368a211e2e89a8efe2dec86cda3647ae15.
\end{center}
We analyze the ordinary
signature Miraidon-S, specifically Algorithms~3--6.  We do not claim an
attack on the separate ring-signature or linkable-ring-signature security
properties.

Let $q$ be a prime, let $M_1,\ldots,M_k\in\F_q^{m\times n}$ be the public
matrices, and let $r$ be the target rank.  For $z\in\F_q^k$, write
\[
  M(z)=\sum_{j=1}^{k} z_jM_j.
\]
Problem~1 in the target asks for $a\in\F_q^k$ such that
$\rank(M(a))\le r$.  It does not state $a\ne\vectzero$.  Thus its literal
homogeneous formulation is satisfied by $a=\vectzero$ for every instance.
Algorithms~3 and~4 instead describe the private witness as an $a$ for which
$E=M(a)$ has rank exactly $r$.  Algorithm~5 generates a nonzero honest secret
and rejects unless the resulting relation has rank $r$; those key-generation
properties do not constrain a malicious signature response.

\subsection{One protocol round}

For an honest round, the prover samples $u\in\F_q^k$ and
$S\in\GL_m(\F_q)$, $T\in\GL_n(\F_q)$, then computes
\begin{align*}
  X &= \sum_{j=1}^{k}u_jSM_jT = S M(u)T,\\
  Q &= SET = S M(a)T,\\
  Q &= LR,
\end{align*}
where $L\in\F_q^{m\times r}$ and $R\in\F_q^{r\times n}$.  It commits to
$c_0=\Com(X,Q)$ and $c_1=\Com(S,T)$.  After a first challenge
$\xi\in\F_q^*$, it sets
\[
  v=u+\xi a,
  \qquad
  y=\Com\!\left(\sum_{j=1}^{k}v_jSM_jT\right).
\]
The verifier then gives a bit challenge $b$:
\begin{itemize}
  \item if $b=0$, the prover opens $X,L,R$; the verifier forms $Q=LR$ and
  checks $c_0$ and $y=\Com(X+\xi Q)$;
  \item if $b=1$, the prover opens a seed for $S,T$ and the vector $v$; the
  verifier checks $c_1$ and
  \[
    y=\Com\!\left(\sum_jv_jSM_jT\right).
  \]
\end{itemize}
Miraidon-S repeats this protocol $t$ times.  The first Fiat--Shamir hash
$\FSone$ produces $t$ nonzero field elements.  The second hash $\FStwo$
produces a bit vector of fixed weight $w$; its support $I$ selects the
$b=1$ rounds.  A seed tree compresses the disclosed seeds, and a Merkle tree
authenticates the branch-zero commitments.

\subsection{Claimed concrete security}

The target states that Miraidon-S is secure against forgery attacks based on
the hardness of MinRank and selects parameters for NIST Levels I, III, and V.
Table~\ref{tab:claimed} records the exact suggested rows used in our tests.
The final column is the target's claimed work exponent for its analyzed
fixed-weight forgery attack.  The target also reports minimum generic MinRank
estimates of approximately $2^{136}$, $2^{200}$, and $2^{265}$ operations for
these rows.  These estimates do not apply to the verifier-acceptance flaw.

\begin{table}[ht]
\centering
\caption{Published Miraidon-S rows and claimed fixed-weight forgery cost.}
\label{tab:claimed}
\small
\begin{tabular}{@{}lrrrrrrrr@{}}
\toprule
Level & $q$ & $m$ & $n$ & $k$ & $r$ & $t$ & $w$ & $\log_2$ cost \\
\midrule
I   & 127 & 10 & 10 & 44 & 4 & 180 & 121 & 128.29 \\
III & 257 & 12 & 12 & 45 & 6 & 317 & 248 & 192.05 \\
V   & 509 & 13 & 13 & 50 & 7 & 512 & 435 & 256.14 \\
\bottomrule
\end{tabular}
\end{table}

\subsection{The Algorithm-6 aggregate-check mismatch}
\label{sec:alg6}

Literal Algorithm~6 is inconsistent with signing.  Algorithm~4 makes $c_0$
the root of a Merkle tree over the per-round $c_{0,i}$ values and makes $c_1$
a flat commitment to all $c_{1,i}$.  The printed verifier applies a flat
commitment to the $c_{0,i}$ values and applies the Merkle verification to the
$c_{1,i}$ values.  The declared signature contains a proof from the $c_0$
tree, so the literal operations do not verify an honestly generated
signature.

Our analysis and implementations use the evident consistency correction
required by the signing algorithm and the surrounding optimization text:
\begin{equation}
  c_0 = \mathsf{VerifyTree}(\{c_{0,i}\}_{i\notin I},\mathsf{Proof}),
  \qquad
  c_1 = \Com(c_{1,1},\ldots,c_{1,t}).
  \label{eq:aggregatefix}
\end{equation}
Honest controls pass this corrected verifier.  The attack below passes the
same verifier.  It is therefore independent of the aggregate-check
transcription error, but the scope of our conclusion is necessarily the
intended, executable verifier rather than literal Algorithm~6.

\section{The zero-witness forgery}

\subsection{The missing exact-rank constraint}

The factor dimensions in an opened branch-zero response imply only
\[
  \rank(LR)\le r.
\]
They do not imply equality, and Algorithm~3 or~6 performs no rank check on
$L$, $R$, or $Q=LR$.  This distinction is load-bearing.  The target's
soundness discussion says that $Q$ has rank $r$ ``by construction,'' but that
statement holds only for an honest prover.  In contrast, the completeness
proof itself describes $Q$ as having rank at most $r$.  Verification never
enforces the exact-rank premise used by extraction.

The most direct malicious choice is
\begin{equation}
  a^*=\vectzero,
  \qquad E^*=Q=L=R=\matzero.
  \label{eq:zeros}
\end{equation}
For any $u\in\F_q^k$, any valid invertible $S,T$, and any
$\xi\in\F_q^*$, set
\[
  v=u,
  \qquad
  X=\sum_{j=1}^{k}u_jSM_jT.
\]
Then the two verifier branches compute the same matrix:
\begin{equation}
  X+\xi Q = X
  =\sum_{j=1}^{k}v_jSM_jT.
  \label{eq:identity}
\end{equation}
The specialization $u=v=\vectzero$ also makes $X=\matzero$.  It removes even
the public-matrix arithmetic: the attacker need not read the matrices
$M_j$.

\begin{theorem}[Universal accepting transcript]
For every Miraidon-S public key, every message, every vector of nonzero first
challenges, and every fixed-weight second challenge, the choices in
Equation~\eqref{eq:zeros} with $u=v=\vectzero$ produce responses accepted by
the intended verifier in every round.
\end{theorem}

\begin{proof}
Before either challenge, the attacker can commit to $X=Q=\matzero$ and to
honestly expanded invertible $S,T$.  For $b=0$, the verifier reconstructs
$Q=LR=\matzero$ and hashes
$X+\xi Q=\matzero$.  For $b=1$, it reconstructs the committed $S,T$ and hashes
$\sum_jv_jSM_jT=\matzero$.  Hence the same delayed commitment
$y=\Com(\matzero)$ is correct for either bit and for every $\xi$.  The attacker
knows all commitment openings, seeds, and authentication paths, so the
corrected aggregate checks in Equation~\eqref{eq:aggregatefix} also pass.
This holds independently in all $t$ rounds.
\end{proof}

\subsection{Forgery algorithm}

The following algorithm spells out the complete noninteractive forgery.  It
uses the target's abstract seed-tree, Merkle-tree, commitment, and
Fiat--Shamir interfaces; concrete encodings are discussed in
Section~\ref{sec:repro}.

\medskip
\noindent\fbox{\parbox{0.94\textwidth}{
\textbf{Algorithm: universal zero-witness forgery.}

\textbf{Input:} a public key $pk$, parameter row
$(q,m,n,k,r,t,w)$, and an attacker-chosen message $\mathsf{msg}$.
\begin{enumerate}
  \item Sample the master seed and salt and derive the $t$ per-round seeds as
  an honest signer would.
  \item For each $i\in[t]$, expand the seed to valid
  $S_i\in\GL_m(\F_q)$ and $T_i\in\GL_n(\F_q)$, and set
  \[
    u_i=v_i=\vectzero,\qquad X_i=Q_i=L_i=R_i=\matzero.
  \]
  Compute
  $c_{0,i}=\Com(X_i,Q_i,\mathsf{Salt},i)$ and
  $c_{1,i}=\Com(S_i,T_i,\mathsf{Salt},i)$.
  \item Build the $c_0$ Merkle tree, let $c_0$ be its root, and set
  $c_1=\Com(c_{1,1},\ldots,c_{1,t})$.
  \item Query
  $(\xi_1,\ldots,\xi_t)=
  \FSone(c_0,c_1,\mathsf{msg},\mathsf{Salt})$.  For every round set
  $y_i=\Com(\matzero)$ and $Y=\Com(y_1,\ldots,y_t)$.
  \item Query
  $b=\FStwo(c_0,c_1,\mathsf{msg},\mathsf{Salt},
  (\xi_i)_{i\in[t]},Y)$ and let $I=\operatorname{supp}(b)$.
  \item Form the honest seed opening for $I$ and the honest Merkle multiproof
  for the opened $c_{0,i}$ values.  For $i\notin I$, output
  $(c_{1,i},X_i,L_i,R_i)$.  For $i\in I$, output $v_i$; the seed path reveals
  the corresponding $S_i,T_i$.
  \item Return the resulting signature
  $(c_0,c_1,\mathsf{Salt},Y,\mathsf{Path},\mathsf{Proof},
  \{\mathsf{resp}_i\}_{i\in[t]})$.
\end{enumerate}
}}
\medskip

The attacker never calls a signing oracle and never evaluates a secret-key
operation.  Because every response is prepared honestly relative to the
malicious zero statement, there is no need to grind either Fiat--Shamir hash.
Binding of the commitment and collision resistance of the Merkle tree do not
help: the attacker opens exactly what it committed.

\section{Complexity and success probability}

The attack performs one ordinary transcript-construction pass.  In the
all-zero specialization it avoids the $k$ public-matrix linear combination
and the rank-$r$ factorization in every round.  Its remaining work consists of
seed expansion, generation of invertible masking matrices, hashing, tree
construction, and serialization.  With the dimensions treated explicitly,
this is polynomial in the parameter sizes and linear in the number of rounds
apart from the cost of producing $S_i,T_i$.  Its signature size is the honest
signature size for the same row.

The success probability is $1$ for every well-formed output of the specified
Fiat--Shamir functions.  In particular, it is independent of $q$, $t$, and
$w$.  This contrasts with the target's $2^{128.29}$, $2^{192.05}$, and
$2^{256.14}$ estimates for challenge-guessing forgeries.  The attack also
does not solve a MinRank instance, so generic MinRank cost estimates are
irrelevant.

Every published row is affected.  The factor matrices with the required
shapes exist as zero matrices for $r=4,6,7$, and both second-challenge branches
occur in the fixed-weight signatures.  The exact branch counts observed in a
deterministic reproduction are given in Table~\ref{tab:allrows}.

\begin{table}[ht]
\centering
\caption{Full-scale all-row reproduction.  ``Fresh'' counts separately
generated attacker-chosen messages.}
\label{tab:allrows}
\small
\begin{tabular}{@{}lrrrccccc@{}}
\toprule
Level & $b=0$ & $b=1$ & time (s) & fresh 1 & fresh 2 & changed & rank fix & honest \\
\midrule
I   & 59 & 121 & 2.684  & pass & pass & reject & reject & pass \\
III & 69 & 248 & 7.329  & pass & pass & reject & reject & pass \\
V   & 77 & 435 & 17.256 & pass & pass & reject & reject & pass \\
\bottomrule
\end{tabular}
\end{table}

Here ``rank fix'' is the exact-rank-hardened verifier applied to the forgery;
``honest'' is an honest signature accepted by that hardened verifier.  All
first challenges were nonzero.  Every opened forged $Q$ had rank zero, while
the honest controls opened matrices of ranks $4$, $6$, and $7$ respectively.

\section{Experimental validation}

\subsection{Primary exact-key experiment}

The primary dependency-free Python implementation instantiated the natural
Algorithm-5 generator at the exact Level-I row
$(127,10,10,44,4,180,121)$.  With deterministic test randomness, key
generation solved the 44-equation/43-ratio system after 199 attempts,
reconstructed the specified 16-element compressed public-key tail, and
confirmed that the honest secret combination had rank~4.

On that key, the same corrected verifier accepted an honest signature and a
forgery on a distinct attacker-chosen message.  The attacker function received
only the public matrices, the fresh message, and attacker randomness.  It did
not receive the secret key and made zero signing queries.  Replaying the
signature on a changed message failed.  Adding an exact-rank test rejected the
forgery and continued to accept the honest signature.

The forge phase took 0.490 seconds.  The complete key-generation,
honest-signing, forging, verification, and control run took 4.94 seconds and
22,528 KiB peak resident memory on one worker.  The recorded environment used
Python~3.12.3 on an AMD EPYC processor under KVM.  This experiment made one
forgery attempt and obtained one acceptance; the algebraic proof above gives
the success statement beyond that single timing trial.

\subsection{Independent implementations}

A fresh-context standard-library implementation independently reconstructed
the complete signature system: finite-field and matrix arithmetic, valid
target-shaped keys, seed expansion, compressed seed disclosure, Merkle
multiproofs, both Fiat--Shamir stages, an honest signer, the attacker, and the
corrected verifier.  It generated two accepted fresh forgeries at each of the
three rows, for an observed $6/6$ full-signature acceptance rate.  All
changed-message replays failed; all exact-rank controls rejected the
forgeries and accepted the honest signatures.  The all-row run took 27.38
seconds and 35,328 KiB peak resident memory on one worker.

A second independent Level-I implementation serialized a public key, an
honest signature, two fresh zero-witness forgeries, a nonzero-$v$ variant, and
a nonzero-factor variant.  It exercised message, response, Merkle-proof,
seed-opening, and $Y$ tampering controls; each malformed signature was
rejected.  Its main run took approximately 5.2 seconds and 24 MiB peak memory.

Finally, a branch probe imported the pinned authors' interactive
proof-of-concept rather than reimplementing its verifier.  At the exact
Level-I dimensions, the zero transcript passed both $b=0$ and $b=1$ for all
126 values of $\xi\in\F_{127}^*$: 252 accepted branch cases out of 252.
Changing an opened $X$ or $v$ while retaining the commitment caused rejection.
This confirms that acceptance follows from the protocol equation rather than
from disabled commitment checks.

\subsection{Controls and repair boundary}

The independent controls delimit what must be checked.
\begin{itemize}
  \item Requiring $v\ne\vectzero$ is insufficient.  Set
  $u=v=e_1$, $Q=0$, and $X=SM_1T$; Equation~\eqref{eq:identity} still holds.
  \item Requiring merely $L\ne0$ and $R\ne0$ is insufficient.  For the
  published $r\ge2$, nonzero factors may be supported on different latent
  coordinates so that $LR=0$.  This variant was accepted.
  \item Requiring full-rank $S$ and $T$ is insufficient; the target already
  does so, and all experiments used invertible masks.
  \item Checking $\rank(Q)=r$ rejects the demonstrated family and preserves
  honest signatures in every tested row.
\end{itemize}

\section{Repair}

The direct verifier repair is to reject every opened branch-zero response
unless
\begin{equation}
  \rank(Q)=\rank(LR)=r.
  \label{eq:rankrepair}
\end{equation}
Equivalently, the verifier may require $L\in\F_q^{m\times r}$ to have full
column rank and $R\in\F_q^{r\times n}$ to have full row rank.  These conditions
force $\rank(LR)=r$ by Sylvester's rank inequality.  Checking $Q\ne0$ alone
blocks this particular zero-$Q$ family, but it does not enforce the exact-rank
relation stated by Algorithms~3 and~4.

A complete specification repair should also:
\begin{enumerate}
  \item define the homogeneous MinRank witness as nonzero or projectively
  normalized and state unambiguously whether the relation is exact rank $r$;
  \item add Equation~\eqref{eq:rankrepair}, or the equivalent two factor-rank
  checks, to Algorithms~3 and~6;
  \item correct the swapped $c_0$ Merkle and $c_1$ flat aggregate checks as in
  Equation~\eqref{eq:aggregatefix}; and
  \item revise the soundness argument so that exact rank is established by a
  verifier check rather than asserted ``by construction.''
\end{enumerate}
The repaired construction requires a fresh proof review; our experiments
show only that the exact-rank check blocks this attack while preserving the
tested honest executions.

\section{Reproducibility and limitations}
\label{sec:repro}

The public reproducer uses only Python's standard library.  From the root of
the public audit repository, run
\begin{verbatim}
python3 attacks/miraidon-s/reproduce.py \
  --output attacks/miraidon-s/result.json
\end{verbatim}
It generates an exact Algorithm-5 Level-I key, checks the honest relation,
signs and verifies an honest message, forges a distinct fresh message without
passing secret material to the attacker, rejects a changed-message replay,
and exercises the exact-rank repair.  The checked-in
\path{attacks/miraidon-s/reference-output.json} records deterministic public
key, message, and signature digests; timing fields are machine-dependent.

The target paper specifies the signature at the algorithm level but does not
provide a complete Sign/Verify implementation, byte encoding, concrete
Fiat--Shamir expansion, or tree encoding.  The reproductions therefore use
unambiguous, domain-separated SHAKE256-based reference encodings and ordinary
binary tree constructions.  This is a limitation for byte-for-byte
interoperability, not for the algebraic result: both branches receive the
same matrix before hashing, and the attacker knows every tree and commitment
opening.  Two implementations with different framing and tree code reached
the same result, and the interactive branch identity passed in the pinned
source.

The experimental timings are reproducibility measurements rather than
optimized benchmarks.  The exact Algorithm-5 run used one deterministic key
and one timed forgery; the independent implementation covered two fresh
messages per parameter row.  The universal success claim follows from
Equation~\eqref{eq:identity}, not statistical extrapolation.  This paper
analyzes version~4 and the intended aggregate-check correction described in
Section~\ref{sec:alg6}; a later specification may differ.  No secret-key
recovery is claimed because the stronger direct-forgery outcome already
breaks message authenticity.

\section{AI and verification disclosure}

This report is an AI-assisted research artifact.  The participation record
is:
\begin{center}
\small
\begin{tabular}{@{}p{0.35\textwidth}p{0.54\textwidth}@{}}
\toprule
Activity & Status \\
\midrule
Target discovery and triage & AI-assisted \\
Attack-history search & AI-run with cached public-source evidence \\
Cryptanalysis & Codex and specialist AI agents \\
Implementation and experiments & AI-generated and AI-run \\
Adversarial review & Separate AI agent \\
Independent reproduction & Fresh-context AI agents and independent code \\
Independent human verification & Pending \\
\bottomrule
\end{tabular}
\end{center}
AI reproduction is not human verification.  The result should be cited and
reviewed with that distinction intact.

\section{Conclusion}

Miraidon-S version~4 accepts a rank-zero factorization in a protocol whose
soundness argument requires exact rank $r$.  The all-zero coefficient vector
therefore gives the forger one prechallenge transcript that answers every
first challenge and both second-challenge branches.  Parallel repetition,
fixed-weight challenges, Fiat--Shamir hashing, and tree compression preserve
rather than remove this identity.  The resulting attack is a practical,
public-key-only fresh-message forgery against every published parameter row.
An exact-rank verifier check blocks the demonstrated forgery, but the relation,
verification algorithm, and soundness proof should be repaired together.

\begin{thebibliography}{10}

\bibitem{miraidon}
R.~Cartor and F.~Slaughter.
\newblock Miraidon: MinRank Identification.
\newblock Cryptology ePrint Archive, Paper 2026/997, version 4, 2026.
\newblock \url{https://eprint.iacr.org/2026/997}.

\bibitem{bardet2020}
M.~Bardet, M.~Bros, D.~Cabarcas, P.~Gaborit, R.~Perlner, D.~Smith-Tone,
  J.-P.~Tillich, and J.~Verbel.
\newblock Improvements of algebraic attacks for solving the rank decoding and
  MinRank problems.
\newblock In \emph{ASIACRYPT 2020}, pages 507--536. Springer, 2020.

\bibitem{bardet2023}
M.~Bardet, P.~Briaud, M.~Bros, P.~Gaborit, and J.-P.~Tillich.
\newblock Revisiting algebraic attacks on MinRank and on the rank decoding
  problem.
\newblock \emph{Designs, Codes and Cryptography}, 91(11):3671--3707, 2023.
\newblock \url{https://doi.org/10.1007/s10623-023-01265-x}.

\bibitem{kales2020}
D.~Kales and G.~Zaverucha.
\newblock An attack on some signature schemes constructed from five-pass
  identification schemes.
\newblock In \emph{CANS 2020}, pages 3--22. Springer, 2020.
\newblock \url{https://doi.org/10.1007/978-3-030-65411-5_1}.

\bibitem{battagliola2025revision}
M.~Battagliola, R.~Longo, F.~Pintore, E.~Signorini, and G.~Tognolini.
\newblock A revision of CROSS security: Proofs and attacks for multi-round
  Fiat--Shamir signatures.
\newblock Cryptology ePrint Archive, Paper 2025/127, 2025.
\newblock \url{https://eprint.iacr.org/2025/127}.

\bibitem{battagliola2025fixed}
M.~Battagliola, R.~Longo, F.~Pintore, E.~Signorini, and G.~Tognolini.
\newblock Security of fixed-weight repetitions of special-sound multi-round
  interactive proofs.
\newblock \emph{Designs, Codes and Cryptography}, 93(10):4323--4353, 2025.
\newblock \url{https://doi.org/10.1007/s10623-025-01686-w}.

\bibitem{courtois2001}
N.~T. Courtois.
\newblock Efficient zero-knowledge authentication based on a linear algebra
  problem MinRank.
\newblock In \emph{ASIACRYPT 2001}, pages 402--421. Springer, 2001.

\bibitem{fiatshamir}
A.~Fiat and A.~Shamir.
\newblock How to prove yourself: Practical solutions to identification and
  signature problems.
\newblock In \emph{CRYPTO '86}, pages 186--194. Springer, 1987.

\bibitem{descala2025}
A.~J. De~Scala and C.~Sanna.
\newblock Smaller public keys for MinRank-based schemes.
\newblock \emph{Journal of Mathematical Cryptology}, 19(1):20240008, 2025.
\newblock \url{https://doi.org/10.1515/jmc-2024-0008}.

\end{thebibliography}

\end{document}
